%! AP Statistics — Unit 4, Lesson 4.12 — Homework
\documentclass[10pt]{article}
\usepackage{apstats-article}
\usepackage{apstats-boxes}

\begin{document}

\pageheader{Unit 4, Lesson 4.12}{Homework}
\namedateperiod

\vspace{0.1in}

\begin{practicebox}
\small
\textbf{Directions:} For each problem, show all formula and calculator setup.
Write complete-sentence interpretations wherever asked.

\vspace{10pt}
\textbf{1.} A standard die is rolled until a 1 appears.
Let $X =$ the number of rolls needed.
\begin{enumerate}[label=(\alph*), leftmargin=2em, itemsep=8pt]
  \item Check BITS and state the distribution.\\
        \writeline\\[1pt]\writeline\\[1pt]\writeline
  \item Compute $P(X = 2)$ and $P(X = 6)$ using the formula. Show setup.\\
        \writeline\\[1pt]\writeline
  \item Find $\mu_X$ and interpret in context.\\
        \writeline\\[1pt]\writeline
  \item Find $P(X > 6) = (1 - p)^6$. Interpret in one sentence.\\
        \writeline\\[1pt]\writeline
\end{enumerate}

\vspace{8pt}
\textbf{2.} A factory produces bolts. Each bolt independently has a 4\% chance
of being too large. Bolts are tested one at a time until the first oversized bolt
is found.
\begin{enumerate}[label=(\alph*), leftmargin=2em, itemsep=8pt]
  \item State the distribution of $X$ (the number of bolts tested).\\
        \blank{5cm}
  \item Compute $P(X = 10)$. Write the formula and answer.\\
        \writeline\\[1pt]\writeline
  \item Find $\mu_X$ and interpret.\\
        \writeline\\[1pt]\writeline
  \item Compute $P(X \le 5)$ using \texttt{geometcdf}.\\
        \writeline
\end{enumerate}

\vspace{8pt}
\textbf{3.} (Binomial vs.\ Geometric) A researcher samples voters.
Each voter independently supports a candidate with probability 0.55.
\begin{enumerate}[label=(\alph*), leftmargin=2em, itemsep=8pt]
  \item The researcher interviews voters until finding the first supporter.
        Let $W =$ the number interviewed. State the distribution of $W$.\\
        \blank{5cm} \quad $\mu_W = $ \blank{3cm}
  \item The researcher interviews exactly 10 voters.
        Let $V =$ the number of supporters. State the distribution of $V$.\\
        \blank{5cm} \quad $\mu_V = $ \blank{3cm}
  \item Explain the key difference between $W$ and $V$ in one or two sentences.\\
        \writeline\\[1pt]\writeline
\end{enumerate}

\vspace{8pt}
\textbf{4.} (AP-Style Free Response) \textit{A certain allele appears in 20\% of
a large population's individuals. A geneticist examines individuals one at a time
(independently) until finding the first one carrying the allele.
Let $X =$ the number of individuals examined.}
\begin{enumerate}[label=(\alph*), leftmargin=2em, itemsep=8pt]
  \item Define the distribution of $X$. Check BITS briefly.\\
        \writeline\\[1pt]\writeline\\[1pt]\writeline
  \item Compute $P(X = 4)$. Show formula setup.\\
        \writeline\\[1pt]\writeline
  \item Find $\mu_X$ and interpret in context.\\
        \writeline\\[1pt]\writeline
  \item Find $P(X \le 3)$. Write the calculator expression and the answer.\\
        \writeline
  \item Find $P(X > 5)$. Interpret in context.\\
        \writeline\\[1pt]\writeline
\end{enumerate}

\vspace{8pt}
\textbf{5. Extension:} The geometric distribution is \textbf{memoryless}.
Verify numerically that $P(X > 5 \mid X > 2) = P(X > 3)$ when $p = 0.20$.
Show both calculations and explain what ``memoryless'' means in the context
of the geneticist problem.

\vspace{4pt}
\writeline\\[1pt]\writeline\\[1pt]\writeline

\end{practicebox}

\end{document}
